% !PRINTED 69
\weberpage{69}
%<<CONT

is to be \(\Delta\). The quantity \(\Delta\) is called the \gesperrt{determinant} of the \(n^2\) elements \(a_i^{(k)}\), more precisely, when the distinction is necessary, an \gesperrt{\(n\)-rowed determinant} (also a determinant of the \(n\)-th degree or \(n\)-th order). The term \(M\) of this sum, that is, the product of all elements standing in the diagonal of the square, is called the \gesperrt{principal term}.

For example, if \(n=2\), we obtain
\begin{equation}
\Delta=a_1^{(1)}a_2^{(2)}-a_2^{(1)}a_1^{(2)},
\tag{4}
\end{equation}
and for \(n=3\) [according to (3) of the preceding paragraph]:
\begin{equation}
\begin{aligned}
\Delta={}&a_1^{(1)}a_2^{(2)}a_3^{(3)}
+a_2^{(1)}a_3^{(2)}a_1^{(3)}
+a_3^{(1)}a_1^{(2)}a_2^{(3)}\\
&-a_1^{(1)}a_3^{(2)}a_2^{(3)}
-a_2^{(1)}a_1^{(2)}a_3^{(3)}
-a_3^{(1)}a_2^{(2)}a_1^{(3)}.
\end{aligned}
\tag{5}
\end{equation}
Or in another notation:
\begin{equation}
\begin{vmatrix}
a&b\\
c&d
\end{vmatrix}
=ad-bc,
\tag{6}
\end{equation}
\begin{equation}
\begin{vmatrix}
a&b&c\\
a'&b'&c'\\
a''&b''&c''
\end{vmatrix}
=ab'c''+bc'a''+ca'b''-ac'b''-ba'c''-cb'a''.
\tag{7}
\end{equation}

The reader is advised to practise the calculation of such determinants on numerical examples.

The notation (1) is in many cases too cumbersome; other, shorter symbols are therefore also in use. Thus \gesperrt{Jacobi}, writing out only the principal term of the expanded determinant, puts
\begin{equation}
\Delta=\sum \pm a_1^{(1)}a_2^{(2)}\cdots a_n^{(n)},
\tag{8}
\end{equation}
and \gesperrt{Kronecker}, still more briefly,
\begin{equation}
\Delta=\left|a_i^{(k)}\right|.
\tag{9}
\end{equation}
Both notations are entirely clear only when the elements are denoted by two indices in the manner presupposed here, and are wholly inapplicable when, for example, the elements are given numerically.

Determinants sometimes occur for which
\[
a_i^{(k)}=a_k^{(i)},
\]
so that in (1) the elements standing symmetrically with respect to the diagonal of the square are equal to one another. In these cases we shall ordinarily place both indices below (to indicate their equivalence), and therefore put
\[
a_{i,k}=a_{k,i}.
\]
Such determinants are called \gesperrt{symmetric}.
